\DocumentMetadata{lang=en, pdfstandard=ua-2, testphase={phase-III,math,table,title,graphic}}
\documentclass[11pt]{article}
\usepackage[top=1.25cm,bottom=1.6cm,left=2.35cm,right=2.35cm]{geometry}
\usepackage[dvipsnames,svgnames,x11names]{xcolor}
\usepackage{amsmath,amssymb,mathtools}
\usepackage{graphicx}
\usepackage{float}
\usepackage{fancyhdr}
\usepackage{hyperref}
\usepackage{titlesec}
\usepackage{tcolorbox}
\usepackage{framed}
\usepackage{listings}
\usepackage{upquote}
\usepackage{array}
\usepackage{pifont}
\hypersetup{pdfauthor={},pdftitle={}}
\setlength{\parindent}{0pt}
\setlength{\parskip}{0.6em}
\pagestyle{fancy}
\fancyhf{}
\renewcommand{\headrulewidth}{0.4pt}
\fancyhead[L]{\fontsize{14}{14}\selectfont EXERCISE}
\fancyhead[R]{\fontsize{14}{14}\selectfont May 11th, 2026}
\setlength{\headheight}{18pt}
\definecolor{codebg}{RGB}{238,238,238}
\definecolor{framegray}{RGB}{140,140,140}
\definecolor{codegreen}{RGB}{80,120,110}
\definecolor{codered}{RGB}{170,70,70}
\definecolor{codeblue}{RGB}{70,70,160}
\definecolor{codepurple}{RGB}{150,60,180}
\definecolor{lightgrayline}{RGB}{210,210,210}
\lstdefinestyle{pythonstyle}{
  backgroundcolor=\color{codebg},
  basicstyle=\ttfamily\small,
  numbers=left,
  numberstyle=\tiny\color{black},
  stepnumber=1,
  numbersep=10pt,
  frame=single,
  rulecolor=\color{lightgrayline},
  framerule=0.8pt,
  xleftmargin=0pt,
  xrightmargin=0pt,
  aboveskip=0.4em,
  belowskip=0.4em,
  showstringspaces=false,
  columns=fullflexible,
  keepspaces=true,
  commentstyle=\itshape\color{codegreen},
  keywordstyle=\color{codeblue},
  stringstyle=\color{codered}
}
\lstdefinestyle{textstyle}{
  backgroundcolor=\color{codebg},
  basicstyle=\ttfamily\small,
  numbers=left,
  numberstyle=\tiny\color{black},
  stepnumber=1,
  numbersep=10pt,
  frame=single,
  rulecolor=\color{lightgrayline},
  framerule=0.8pt,
  xleftmargin=0pt,
  xrightmargin=0pt,
  aboveskip=0.4em,
  belowskip=0.4em,
  showstringspaces=false,
  columns=fullflexible,
  keepspaces=true
}
\titleformat{\section}{\normalfont\bfseries\fontsize{24}{26}\selectfont}{}{0pt}{}
\titleformat{\subsection}{\normalfont\bfseries\fontsize{17}{19}\selectfont}{}{0pt}{}
\newcommand{\problemheader}[1]{%
  \begin{tcolorbox}[colback=gray!10,colframe=gray!55,boxrule=0.7pt,arc=2mm,left=10pt,right=10pt,top=8pt,bottom=8pt]
  {\fontsize{17}{19}\selectfont \textbf{#1}}
  \end{tcolorbox}%
}
\newcommand{\inlinecode}[1]{\colorbox{gray!15}{\texttt{#1}}}
\begin{document}

\begin{figure}[H]
  \centering
  \includegraphics[width=0.62\textwidth]{img-0.jpeg}
  \caption{Left: The gridworld map considered in this exercise. Right: The state-action function $Q_{\pi}(s,a)$ for the random policy computed using policy evaluation.}
\end{figure}

\section{2\quad Computing State-action values using random probe-
droids (\texttt{probe\_droids.py})}

In this problem, you a dispatching a large amount of probe-droids in a meteor field shown in fig. 2 (left). Probe droids fly around randomly\footnote{i.e., their policy $\pi$ selects among all available actions with equal probability} , and by flying a large number $n$ of probe-droids we can approximate the state-action function $Q_{\pi}(s,a)$ and use it for nefarious purposes.

For reference, fig. 2 (right) shows the (true) action-value function computed using policy evaluation $(\gamma = 1)$, however, we have to compute it using samples.

One way to do that is to use standard tabular Sarsa learning, and simulate the random policy by setting $\epsilon = 1$ (thereby guaranteeing random actions). The result of this, using $n = 20000$ episodes and a learning rate of 0.02, can be seen in (left). As can be seen, it fairly closely matches policy evaluation but with some randomness due to the learning rate/finite number of episodes.

\textbf{This exercise:} \quad Your task is to estimate $Q_{\pi}(s,a)$, for the random policy, using two different linear function approximators. This is equivalent to evaluating Sarsa with a (suitable) linear function approximator using $\epsilon = 1$.

We construct the linear function approximator as follows: Suppose the gridworld has a height of $H = 3$ tiles and a width of $W = 3$ tiles.

A state has at most $\mathcal{A}(s)=\{0,1,2,3\}$ actions. Suppose we let $\boldsymbol{x}(s)$ be any function approximator depending on a state $s$, and let $l=|\boldsymbol{x}(s)|$ be the dimension.

We can then define the full linear function approximator as a $4l$ long vector $\boldsymbol{x}(s,a)$ defined as:

\begin{equation*}
\boldsymbol{x}(s,a=0)=
\begin{bmatrix}
\boldsymbol{x}(s)\\
\mathbf{0}_{3l}
\end{bmatrix},\quad
\boldsymbol{x}(s,a=1)=
\begin{bmatrix}
\mathbf{0}_{l}\\
\boldsymbol{x}(s)\\
\mathbf{0}_{2l}
\end{bmatrix},\quad
\boldsymbol{x}(s,a=2)=
\begin{bmatrix}
\mathbf{0}_{2l}\\
\boldsymbol{x}(s)\\
\mathbf{0}_{l}
\end{bmatrix},\quad
\boldsymbol{x}(s,a=3)=
\begin{bmatrix}
\mathbf{0}_{3l}\\
\boldsymbol{x}(s)
\end{bmatrix},
\end{equation*}
where e.g. $\mathbf{0}_{3l}$ is a vector of $3l$ zeros.

Thus, all that remains is to define $\boldsymbol{x}(s)$.

\newpage

\begin{figure}[H]
  \centering
  \includegraphics[width=0.94\textwidth]{img-2.jpeg}
  \caption{Left: The gridworld map considered in this exercise where the $Q(s,a)$ function is estimated using Sarsa larning. Middle: The $Q(s,a)$ function estimated using the (linear) function approximator given below. Right: $Q(s,a)$ estimated using the quadratic function approximator. Note that you will not get these exact results due to the finite number of samples and $\alpha > 0$.}
\end{figure}

\textbf{Testing and evaluation}\quad After you have implemented the feature approximator, your task is to evaluate select values of $Q(s,a)$ for specific state/action combinations. For example, the following code shows how this can be done using the Sarsa agent:

\begin{lstlisting}[style=pythonstyle,language=Python]
# probe_droids.py
agent = SarsaAgent(env, alpha=0.02, gamma=1., epsilon=1)
train(env, agent, num_episodes=10000, verbose=False)

states_and_actions = [((0, 0), 0),
                      ((0, 0), 1),
                      ((2, 0), 3),
                      ((0, 2), 1),
                      ]
q_values = {(s, a): agent.Q[s, a] for s, a in states_and_actions}
for (s, a), q in q_values.items():
    print(f"In state {s=} and action {a=} we have Q(s,a) = {q}.")
\end{lstlisting}

This will produce output which is consistent with figure fig. 2, right:

\begin{lstlisting}[style=textstyle]
In state s=(0, 0) and action a=0 we have Q(s,a) = 0.39042215523071055.
In state s=(0, 0) and action a=1 we have Q(s,a) = -0.30793653501737517.
In state s=(2, 0) and action a=3 we have Q(s,a) = -0.34343217577946533.
In state s=(0, 2) and action a=1 we have Q(s,a) = 0.9999999999999947.
\end{lstlisting}

\textbf{The linear function approximator}\quad For the first feature map, we use the normalized grid coordinates of the state. Let $s=(i,j)$ denote the row and column of the agent, and

\newpage

define

\begin{equation}
\boldsymbol{z}_{\mathrm{lin}}(s)=
\begin{bmatrix}
1\\
u\\
v
\end{bmatrix},
\qquad
u=\frac{i}{H-1}\text{ and }v=\frac{j}{W-1}.
\tag{2}
\end{equation}

Thus, the state feature dimension is $l=3$.

\begin{tcolorbox}[colback=gray!10,colframe=gray!55,boxrule=0.7pt,arc=2mm,left=12pt,right=12pt,top=8pt,bottom=10pt]
{\fontsize{16}{18}\selectfont Problem 2 \textit{Linear function approximator}}\par
\vspace{0.35em}
Implement the function \inlinecode{linear\_experiment}. The function should accept a number of episodes, the learning rate $\alpha$, a gridworld layout, and a sequence of states and actions. It should then use the linear function approximator described in eq. (2) together with the action-dependent construction above, and return a dictionary of q-values in the previously described format. An example of what you can expect is shown in fig. 3.
\end{tcolorbox}

When done, the script will print out the $Q$-values like this (note: these values are somewhat random due to the limited number of samples)

\begin{lstlisting}[style=textstyle]
In state s=(0, 0) and action a=0 we have Q(s,a)=0.2226488621622873 (Linear function approximator)
In state s=(0, 0) and action a=1 we have Q(s,a)=-0.19386140154078677 (Linear function approximator)
In state s=(2, 0) and action a=3 we have Q(s,a)=-0.10626710421845728 (Linear function approximator)
In state s=(0, 2) and action a=1 we have Q(s,a)=0.46967508585248996 (Linear function approximator)
\end{lstlisting}

\begin{minipage}[t]{0.06\textwidth}
\vspace{0pt}
{\Large\ding{108}}
\end{minipage}
\begin{minipage}[t]{0.9\textwidth}
\vspace{0pt}
\textbf{Info:}
\begin{itemize}
\item One approach is to use the \inlinecode{LinearSemiGradSarsa}, using $\epsilon = 1$, but specify a different feature encoder
\item you can find more about feature encoders here: \url{https://www2.compute.dtu.dk/courses/02465/exercises/ex11.html}
\item You can get the $H$ and $W$ values using \inlinecode{env.mdp.height} and \inlinecode{env.mdp.width}.
\end{itemize}
\end{minipage}

\textbf{The quadratic function approximator}\quad For the second feature map, we augment the normalized coordinates with quadratic terms and an interaction term. Using the same notation as above, define

\begin{equation}
\boldsymbol{z}_{\mathrm{quad}}(s)=
\begin{bmatrix}
1\\
u\\
v\\
u^{2}\\
v^{2}\\
uv
\end{bmatrix},
\qquad
u=\frac{i}{H-1}\text{ and }v=\frac{j}{W-1}.
\tag{3}
\end{equation}

In this case, the state feature dimension is $l=6$.

\newpage
\thispagestyle{empty}

\noindent
\begin{minipage}[t]{\textwidth}
  \vspace*{-0.2cm}
  \begin{minipage}[t]{0.5\textwidth}
    \raggedright
    {\large EXERCISE}
  \end{minipage}%
  \begin{minipage}[t]{0.5\textwidth}
    \raggedleft
    {\large May 11th, 2026}
  \end{minipage}
\end{minipage}

\vspace{0.15cm}
\noindent\rule{\textwidth}{0.4pt}

\vspace{1.2cm}

\begin{tcolorbox}[enhanced,colback=white,colframe=gray!70,boxrule=0.4pt,arc=2.5mm,left=10pt,right=10pt,top=12pt,bottom=12pt]
\noindent
\raisebox{0.8ex}{\hspace{0.4em}\fboxsep=3pt\colorbox{white}{\strut \large Problem 3 \textit{Quadratic function approximator}}}

\vspace{0.45em}

Implement the function \texttt{quadratic\_experiment}. The function should accept a number of episodes, the learning rate $\alpha$, a gridworld layout, and a sequence of states and actions. It should then use the quadratic function approximator described in eq.\ (3) together with the action-dependent construction above, and return a dictionary of q-values in the previously described format. An example of what you can expect is shown in fig.\ 3.
\end{tcolorbox}

\vspace{0.9cm}

\begin{center}
\parbox{0.88\textwidth}{\centering Again, the following are examples of valid $q$-values, but yours may vary somewhat due to the finite number of samples:}
\end{center}

\vspace{0.55cm}

{\footnotesize
\begin{tabular}{@{}r l@{}}
1 & In state s=(0, 0) and action a=0 we have Q(z,a)=0.28231914123335555 (Quadratic function approximator)\\
2 & In state s=(0, 0) and action a=1 we have Q(z,a)=-0.1904661084415149 (Quadratic function approximator)\\
3 & In state s=(2, 0) and action a=3 we have Q(z,a)=-0.17901693611399366 (Quadratic function approximator)\\
4 & In state s=(0, 2) and action a=1 we have Q(z,a)=0.7379605641504839 (Quadratic function approximator)
\end{tabular}
}

\vspace{1.6cm}

{\huge \textbf{References}}

\vspace{0.8cm}

\noindent [SB18]\hspace{0.8em}Richard S. Sutton and Andrew G. Barto. \textit{Reinforcement Learning: An Introduction}.\\
\hspace*{3.9em}The MIT Press, second edition, 2018. (Freely available online).

\end{document}